\chapter{ViFailback Case Study: Application to a Further Dataset}
\label{ch:vifailback}

\section{Motivation and Replication Objectives}
On held-out ViFailback manipulation tasks, the logistic sensor model ranked episode
outcomes more accurately than the fitted combination of video and sensor scores. The
dataset provides three camera views and robot observations for complete episodes
\citep{Zeng2026ViFailback}. The study compared the two branch predictions and their
combination on the same tasks (SQ1 and SQ2). A separate comparison isolated the choice
between a generic frozen visual encoder and one initialized from a REASSEMBLE
checkpoint, with a new prediction head fitted in each case.

The primary replication contrasts original frozen vision, logistic sensing, averaging
and F2. HGB supplied a stronger numerical reference, while camera ablations and donor
transfer tested visual representations. The primary comparison was specified on 30
October 2026 after some early results had been observed. Its interpretation is therefore
partly retrospective.

The reporting cutoff for this revision is extended through the completed-scope analysis
of 24 November 2026. In addition to the frozen and donor comparisons, FULL20 completed
seed~42, outer fold~0 and its three inner fits; FULL100 completed its single-role pilot.
The researcher then stopped further training. The remaining target-CNN matrix and SQ3
corruption study are unfinished, not negative model results. All reported comparisons
concern intact inputs.\ev{Original and extended scope}{\EvidenceRef{E102}{},
\EvidenceRef{E096}{}, \EvidenceRef{E115}{}, \EvidenceRef{E116}{}}

\section{Dataset Overview and Prediction Task}
The native source contains 5,202 episodes with 4,545 failure and 657 no-failure
annotations across 100 canonical tasks and 109 source directories. The evaluation pins
both the original release and its converted LeRobot representation
\citep{ViFailbackHF2025}. Matching version identities matters because the paper, source
directory layout, converted observations and admitted model cohort describe different
levels of the evidence.

The source paper evaluates visual-symbol and visual-language reasoning, including
open-ended failure explanations and policy interventions. Its model does not consume the
280 engineered robot-observation features used here. Its benchmark also differs in task,
cohort and metric: a reported visual-language benchmark score is not failure-positive or
both-class AP for this classifier. The paper therefore provides dataset and
completed-outcome context, rather than a directly comparable fusion score.

The final cohort contains 4,382 episodes, including 3,753 failures and 629 successes
across 99 canonical tasks. Failure is the positive class, giving prevalence 0.856458. An
episode contributes one target, although it supplies 48 sampled RGB images. The task is
retrospective prediction of native whole-episode annotations. It is not frame-level
failure localization, verification of a task's final mechanical success independently of
annotation, or advance warning before a failure occurs.\ev{Native and final ViFailback
labels}{\EvidenceRef{E087}{}}

\section{Dataset Challenges and Cohort Construction}
The provisional core contained 4,384 episodes after 818 unique exclusions or
quarantines. Full three-camera quality checks then identified two defective camera
streams: episode~2071's right wrist and episode~4784's left wrist. Their source bytes
remain preserved. Removing them from analysis produces the final 4,382 common-support
episodes and 820 unique exclusions. Reason counts can overlap; summing them as if
disjoint would overstate the excluded population.

The records distinguish exact trajectory/video duplication and invalid telemetry from
low visual quality and uncertain semantics. Duplication and invalid observations can
invalidate the input contract; visual difficulty does not justify excluding examples
after inspecting correctness. Coordinate task~55 remains quarantined rather than
receiving a guessed sign correction. Agreement of converted 25~Hz observations does not
independently certify original hardware synchronization. Canonical task separation also
leaves operator/session identity and semantic similarity between tasks unresolved,
limiting the independence claim to the grouping actually checked.\ev{Cohort validity and
remaining dependence}{\EvidenceRef{E093}{}}

\section{Preparation and Pre-Checks}
Sixteen uniform, matching frame indices are selected in each of three cameras. These are
sixteen temporal positions and 48 images, not 48 distinct timestamps. The source image
processor resizes to $640\times640$, divides RGB values by 255 and does not apply
ImageNet normalization. Equal camera and temporal means remove order information from
the episode embedding. A fixed crop shown in an example is context visualization, not a
detected object or validated failure location.

The sensor adapter retains fourteen position, twelve nonplaceholder velocity and two
gripper-effort fields. Ten statistics per channel produce 280 values. These observations
are not calibrated wrench or force/torque measurements. Labels, failure categories,
action commands, keyframes and task/source identifiers do not enter ordinary model
features. Per-episode feature calculation is separated from population fitting:
standardizers use T only, and calibration uses disjoint C examples.

The twenty-example visual pilot was selected for observability and balance. Nineteen
decided observations agreed with native labels; B02 remained uncertain. This is a
bounded project examination, not an independent human-rated prevalence sample or a
held-out classifier benchmark. The additional automated consistency checks are also not
a signed second-human review. Technical decoding and a successful native-label join do
not resolve every episode's semantic meaning.\ev{ViFailback
preprocessing}{\EvidenceRef{E088}{}}\ev{Scope of independent
checks}{\EvidenceRef{E092}{}}

\section{Reused Components and Dataset-Specific Adaptations}
The original visual branch loaded the ResNet18-vd CNN backbone inside the pretrained
RT-DETR-R18vd checkpoint \citep{Zhao2024RTDETR}. The extractor stopped at the CNN
feature maps, before the hybrid Transformer encoder, detection decoder, and box heads.
Spatial means of three maps with 128, 256, and 512 channels yielded 896 features per
image. The branch averaged these features across frames and cameras and fitted a new
$896\rightarrow128\rightarrow1$ head to predict the episode outcome. Three-view pooling,
the sensor statistics, and the task partitions were specific to ViFailback.

The logistic sensor branch uses balanced L2 regression with $C=1$, training-only scaling
and a separately fitted sigmoid calibrator. Averaging and F2 consume its calibrated
probabilities alongside the visual branch. F2 is unweighted L2 logistic probability
stacking; HGB is a stronger separate reference and never replaces the sensor input
inside original F2. Consequently, a gain over vision or logistic sensing cannot be
reported as a gain over the strongest available unimodal method.

The donor arms transferred compatible CNN tensors from a completed source checkpoint and
fitted a fresh ViFailback head, scaler, calibrator and fusion model. The FROZEN and
DONOR conditions held the CNN weights fixed. The FULL conditions train the executed CNN
convolutions and normalization affine terms while holding running moments fixed; their
completed coverage is stated below. This is CNN adaptation, not training the detector
encoder, decoder or detection heads. The completed DONOR comparison therefore measures
supervised target-head fitting on a transferred frozen representation.\ev{Matched model
definitions}{\EvidenceRef{E097}{}}

\section{Evaluation Protocol}
The original protocol uses three outer task-disjoint folds and three inner OOF
assessments within each outer-training pool. Each branch separates gradient-fitting T,
checkpoint-selection V, calibration C and assessment A task groups. A final branch is
not automatically fitted to the entire outer-training pool: its T/V/C roles remain
distinct. The inner assessment probabilities build a meta bank without using examples
that fitted, selected or calibrated their generating branch.

Checkpoint selection uses validation macro AP; the selected checkpoint is then
calibrated on C, and assessment uses a fixed threshold of 0.5. Failure AP, success AP
and their arithmetic mean are retained. With failure prevalence above 85\%,
failure-positive AP alone can conceal poor success discrimination. Pooled AP is
calculated from disjoint outer predictions, separately for each seed. It is not the
arithmetic mean of fold APs, and seed repetition does not add independent episodes.

The later matched twenty-epoch arms use the first twenty epochs of a fixed cosine-50
schedule, while the hundred-epoch arms use cosine-100. Highest V macro AP selects the
checkpoint, then lower V BCE and the earliest epoch resolve ties. All configured epochs
are intended to execute even when the chosen checkpoint is earlier. The twenty-epoch
study has a complete three-inner OOF bank. The hundred-epoch pilot instead reserves 788
meta episodes under one predeclared role, followed by 1,107 outer assessments. It is not
a full nested three-fold hundred-epoch experiment.\ev{Original evaluation
contract}{\EvidenceRef{E094}{}}\ev{Matched training contract}{\EvidenceRef{E098}{}}

\section{Experimental Investigations}
The original study evaluated the frozen visual branch, its high-camera control, sensor
models and clean fusion. By 24 November, the matched study had 79 completed fits of 111
planned: 36 each for FROZEN20 and DONOR20, four for FULL20, and one for each
hundred-epoch pilot. These support 22 outer fusion assessments: eighteen frozen/donor
twenty-epoch assessments, one FULL20 assessment and three pilots. Conditions and seeds
reuse episodes. The 32 remaining FULL20 fits contribute no performance
estimates.\ev{Extended matched evaluation scope}{\EvidenceRef{E115}{},
\EvidenceRef{E116}{}}

\begin{table}[htbp]
\caption{Final Scope of the Matched ViFailback Comparisons and Robustness
Tests}\label{tab:vif-status}
\centering\small

\begin{tabularx}{\linewidth}{@{}lX@{}}\toprule
Condition & Evaluation coverage\\\midrule
\multicolumn{2}{@{}l}{\textit{Completed final scope}}\\
FROZEN20 & 36 fits across three seeds and three outer folds\\
DONOR20 & 36 fits across three seeds and three outer folds\\
FULL20 & Four fits: seed 42, outer fold 0 and its three inner fits\\
FROZEN100 / FULL100 / DONOR100 & One single-role pilot fit each\\\addlinespace
\multicolumn{2}{@{}l}{\textit{Unfinished / not evaluated}}\\
Remaining FULL20 & 32 planned fits; researcher stopped further training\\
SQ3 corruption tests & Missing, degraded and temporally shifted inputs\\\bottomrule
\end{tabularx}
\displaynote{Coverage through 24 November 2026. Completed fits and inconclusive
contrasts are distinct from unexecuted conditions. All completed comparisons used intact
inputs.}
\end{table}
\ev{Matched fit coverage}{\EvidenceRef{E115}{}, \EvidenceRef{E116}{}}

\subsection{Unimodal Reference Models}
The original frozen visual head completed all 36 planned fits, as did its
high-camera-only control. The logistic sensor branch is deterministic under the fixed
partitions. S1 and S3 denote the same fitted recipe. HGB supplies a stronger numerical
comparison. Table~\ref{tab:vif-bothclass} gives the pooled seed-42 scores from the
thesis experiment. Reporting average precision for each outcome makes the disparity
between failure and success ranking visible.

\begin{table}[htbp]
\caption{ViFailback Prediction Performance by Outcome Class}\label{tab:vif-bothclass}
\centering\small

\begin{tabular}{lrrrr}\hline Model&Failure AP&Success AP&Macro AP&AUROC\\\hline
Frozen visual V1&0.853052&0.137260&0.495156&0.479402\\
Logistic sensor&0.954370&0.343577&0.648974&0.784370\\
Sensor HGB&0.970510&0.429530&0.700020&0.853756\\
Averaging F1&0.931760&0.289037&0.610399&0.716501\\
Original F2&0.945168&0.312405&0.628787&0.752701\\\hline
\end{tabular}
\displaynote{Scores pool the original seed-42 results separately for each class.}
\end{table}
\ev{Original both-class scores}{\EvidenceRef{E089}{}}

The completed target-CNN comparisons are assessed against their matched frozen controls
below. Their existence does not imply completion of the three-fold, three-seed
adaptation matrix.

\subsection{Original Fusion-Procedure Comparison}
In the original frozen comparison, F2 improves over vision and averaging but has lower
pooled macro AP than the logistic sensor branch for all three seeds: 0.628787, 0.625400
and 0.621067 against 0.648974. The sensor HGB point score is higher still. These results
support a narrower conclusion than ``fusion works'': the numerical information can
improve weak visual prediction, but adding the visual branch does not improve the
strongest component used in F2.

The later matched twenty-epoch comparison also favored the numerical reference.
Table~\ref{tab:vif-matched} lists every seed's pooled score after all three outer
assessments complete. Despite small fold-specific differences around logistic sensing,
FROZEN20's pooled F2 scores are lower by 0.019694, 0.021935 and 0.024419. Saved
conditional task-cluster intervals for these differences exclude zero. This is a
model-and-protocol-specific negative comparison, not proof that every visual
representation will damage fusion.

\begin{table}[htbp]
\caption{ViFailback Performance in the Matched Comparison}\label{tab:vif-matched}
\centering\small

\begin{tabular}{lrrrr}\hline Arm&Seed&Visual&F2&Sensor LR\\\hline
FROZEN20&42&0.507865&0.629280&0.648974\\
FROZEN20&123&0.495229&0.627039&0.648974\\
FROZEN20&20261023&0.487612&0.624555&0.648974\\
DONOR20&42&0.491368&0.627750&0.648974\\
DONOR20&123&0.489011&0.637852&0.648974\\
DONOR20&20261023&0.492185&0.630936&0.648974\\\hline
\end{tabular}
\displaynote{Macro average precision was calculated from pooled predictions over three outer folds.}
\end{table}
\ev{Matched complete outer comparisons}{\EvidenceRef{E096}{}}

\subsection{Representation-Transfer Comparison}
Following donor feature extraction, DONOR20 completed 36 target-head fits. Its source
completed 100 epochs and selected epoch~47; 115 compatible CNN tensors are loaded,
without source head, scaler, calibrator or F2 coefficients. This establishes an executed
learned-representation transfer experiment with fresh target supervision. The fixed
donor is one source-training realization; three target seeds do not replicate donor
selection or training.

It does not establish a consistent predictive advantage. Relative to FROZEN20, donor
visual pooled macro-AP differences are $-0.016497$, $-0.006219$ and $+0.004573$. Their
saved task-cluster intervals cross zero. Donor F2 differences are mixed: $-0.001530$,
$+0.010813$ and $+0.006381$. F2 with the transferred encoder achieved higher macro-AP
than F2 with the frozen encoder for two of the three target seeds. The sensor-only
logistic model achieved higher macro-AP than the transferred-encoder F2 model for all
three seeds.

On the same 1,107 outer episodes, the hundred-epoch frozen-backbone pilot achieved macro
AP of 0.550995 for vision and 0.694219 for F2. This arm reused the completed historical
head-training trajectory after verifying the same checkpoint selection. The
donor-backbone pilot achieved 0.571448 and 0.696600, respectively. Both pilots used a
single seed and fit role, with a separate 788-episode bank for fitting fusion. Neither
pilot updated the target CNN.\ev{Donor identity and pilot fit
provenance}{\EvidenceRef{E098}{}, \EvidenceRef{E099}{}, \EvidenceRef{E100}{}}

\paragraph{Completed target-CNN adaptation.}
On the same 1,107 outer-fold-0 episodes, FULL20 achieved visual macro AP .492370 and F2
.697858; FULL100 achieved .544332 and .699260. The corresponding sensor LR and HGB
references were .707996 and .757878. Table~\ref{tab:vif-adaptation} reports the paired
differences against each arm's own frozen control. Neither visual nor F2 contrast
resolves improvement or deterioration; an interval including zero does not establish
equivalence. The twenty-epoch meta bank contains 3,275 inner-OOF episodes, whereas the
hundred-epoch pilot uses 788 held-out meta episodes. Comparisons between these protocols
cannot isolate training duration.

\begin{table}[htbp]
\caption{Completed ViFailback CNN-Adaptation Contrasts}\label{tab:vif-adaptation}
\centering\small
\begin{tabularx}{\linewidth}{@{}lYY@{}}\toprule
Contrast & Visual macro-AP difference & F2 macro-AP difference\\\midrule
FULL20 $-$ FROZEN20 & $+.004946$ [$-.027460$, $+.040927$] & $-.009260$ [$-.026854$, $+.005016$]\\
FULL100 $-$ FROZEN100 & $-.006662$ [$-.027982$, $+.027078$] & $+.005040$ [$-.019473$,
$+.029255$]\\\bottomrule
\end{tabularx}
\displaynote{Seed 42, outer fold 0 only. Paired 95\% intervals resample canonical tasks,
conditional on fixed models; they omit refitting, selection and unverified
session/operator dependence. Twenty- and hundred-epoch fusion-fitting protocols differ.}
\end{table}
\ev{Completed adaptation results and recorded independent arithmetic}{\EvidenceRef{E115}{}}

\subsection{Scope of the Robustness Evaluation}
ViFailback contributes clean-input comparisons to SQ1 and SQ2. The SQ3 corruption study
remained unexecuted at the extended reporting cutoff. The specified visual interventions
included missing frames or windows, blur, occlusion, camera loss, JPEG degradation and
within-episode timing shifts. Sensor interventions included missing channels or
intervals, Gaussian noise, frozen readings and timing shifts. The matched CNN comparison
contained no quality gate. Performance under these interventions remains a question for
future work.

The specified availability rules route complete single-modality absence to the surviving
branch and abstain when both inputs are absent. These rules do not estimate reliability
from data. The edge-padded sampling shifts simulate timing errors without measuring
drift between hardware clocks.\ev{Specified corruption and fallback
conditions}{\EvidenceRef{E091}{}}

\subsection{Additional Exploratory Baselines}
HGB uses 200 iterations, learning rate 0.05, L2 penalty 0.1 and minimum leaf size
twenty; leaves 7, 15 and 31 are selected on V macro AP, followed by C-only calibration.
Its search is real candidate selection, unlike the fixed MSM recipe or a claim that
every parameter was optimized. It remains secondary and does not change the prespecified
F2 input after its favorable assessment.

The high-camera control fitted a new head on sixteen high-view frames under the original
cohort and split rules. The original frozen study also evaluated a clean
quality-conditioned gate and its availability-only control. Estimating the effect of
corruption exposure would require the deferred augmented and stress
comparisons.\ev{Original controls and fixed recipes}{\EvidenceRef{E094}{},
\EvidenceRef{E090}{}, \EvidenceRef{E091}{}}

\section{Results, Interpretation, and Limitations}
In the original comparison, F2 ranked failures (AP near 0.95) more accurately than
successes (Table~\ref{tab:vif-bothclass}). In outer fold~0 of the matched FROZEN20
comparison, seed~42 gave F2 macro AP of 0.707118 and a Brier score of 0.153928, yet
balanced accuracy was 0.5 at the fixed threshold. Selecting a threshold on those same
assessment episodes would bias the estimate of decision performance.\ev{Saved matched
secondary metrics}{\EvidenceRef{E095}{}}

Some fitted sigmoid calibrators reversed the ranking of episodes within a fit because
their slopes were negative. Even positive slopes can change the ordering of episodes
from different folds when the fitted maps differ. Comparing raw and calibrated rankings
helps interpret differences between pooled AP and mean-fold AP.

The three donor-minus-frozen visual contrasts had task-cluster intervals spanning zero.
These paired intervals used 2,000 draws and held the fitted models, donor and split
assignments fixed. They omit uncertainty from refitting, donor selection and possible
dependence within sessions or operators. The exploratory comparisons had no multiplicity
adjustment. Coefficients were fixed before outer scoring, although the labels and some
tasks had already been examined during development.\ev{Task-cluster uncertainty and
evaluation limits}{\EvidenceRef{E101}{}}

The complete twenty-epoch frozen/donor comparison addresses representation transfer. The
November 24 extension additionally assesses target adaptation on one outer fold, with
inconclusive paired effects. The rest of the adaptation matrix remains unfinished, and
no completed matched corruption comparison establishes robustness.

\section{Relationship to the Earlier Case Studies}
In MSM, the logistic sensor model achieved higher abnormal-cutting AP, averaged over
seed-specific pooled predictions, than the original stacker. In ViFailback, it achieved
higher macro-AP in both the original comparison and the later matched frozen comparison.
In HANDOVER, the logistic sensor model also outperformed averaging. Missing participant
identities precluded a learned-fusion comparison. The direction of the effect can be
compared across these tasks, although their absolute AP values refer to different
outcomes.

Transferring the donor CNN produced mixed effects across target seeds. For complete
ViFailback episodes under the task-disjoint protocol, the logistic sensor predictor
ranked outcomes more accurately than F2 in the matched twenty-epoch comparison. The
completed adaptation contrasts do not establish a change in that conclusion, but cannot
resolve adaptation across the full cohort. Effects of input degradation remain
unmeasured.
